\begin{lemma}
\label{lemma-helper-finite-type}
Let $(A, I)$ be a pair. Let $A \to B$ be a finite type ring map
such that $B/IB = C_1 \times C_2$ with $A/I \to C_1$ finite.
Let $B'$ be the integral closure of $A$ in $B$.
Then we can write $B'/IB' = C_1 \times C'_2$ such that
the map $B'/IB' \to B/IB$ preserves product decompositions
and there exists a $g \in B'$ mapping to $(1, 0)$ in
$C_1 \times C'_2$ with $B'_g \to B_g$ an isomorphism.
\end{lemma}

\begin{proof}
Observe that $A \to B$ is quasi-finite at every prime of the
closed subset $T = \Spec(C_1) \subset \Spec(B)$ (this follows
by looking at fibre rings, see
Algebra, Definition \ref{algebra-definition-quasi-finite}).
Consider the diagram of topological spaces
$$
\xymatrix{
\Spec(B) \ar[rr]_\phi \ar[rd]_\psi & & \Spec(B') \ar[ld]^{\psi'} \\
& \Spec(A)
}
$$
By Algebra, Theorem \ref{algebra-theorem-main-theorem}
for every $\mathfrak p \in T$ there is a $h_\mathfrak p \in B'$,
$h_\mathfrak p \not \in \mathfrak p$ such that $B'_h \to B_h$ is
an isomorphism. The union $U = \bigcup D(h_\mathfrak p)$ gives an open
$U \subset \Spec(B')$ such that $\phi^{-1}(U) \to U$ is a homeomorphism
and $T \subset \phi^{-1}(U)$. Since $T$ is open in $\psi^{-1}(V(I))$
we conclude that $\phi(T)$ is open in $U \cap (\psi')^{-1}(V(I))$.
Thus $\phi(T)$ is open in $(\psi')^{-1}(V(I))$.
On the other hand, since $C_1$ is finite over $A/I$ it is
finite over $B'$. Hence $\phi(T)$ is a closed subset of $\Spec(B')$
by Algebra, Lemmas \ref{algebra-lemma-going-up-closed} and
\ref{algebra-lemma-integral-going-up}. We conclude that
$\Spec(B'/IB') \supset \phi(T)$ is open and closed. By
Algebra, Lemma \ref{algebra-lemma-disjoint-implies-product}
we get a corresponding product decomposition $B'/IB' = C'_1 \times C'_2$.
The map $B'/IB' \to B/IB$ maps $C'_1$ into $C_1$ and $C'_2$ into $C_2$
as one sees by looking at what happens on spectra (hint: the inverse
image of $\phi(T)$ is exactly $T$; some details omitted).
Pick a $g \in B'$ mapping to $(1, 0)$ in $C'_1 \times C'_2$
such that $D(g) \subset U$; this is possible because $\Spec(C'_1)$
and $\Spec(C'_2)$ are disjoint and closed in $\Spec(B')$ and
$\Spec(C'_1)$ is contained in $U$. Then $B'_g \to B_g$ defines a homeomorphism
on spectra and an isomorphism on local rings (by our choice of $U$ above).
Hence it is an isomorphism, as follows for example from
Algebra, Lemma \ref{algebra-lemma-characterize-zero-local}.
Finally, it follows that $C'_1 = C_1$ and the proof is complete.
\end{proof}
